\documentclass{article}
\usepackage{amsmath,amssymb,amsthm}
\usepackage{enumitem}
\newtheorem{theorem}{Theorem}
\newtheorem{lemma}{Lemma}
\newtheorem{assumption}{Assumption}
\newcommand{\prox}{\operatorname{prox}}
\newcommand{\argmin}{\operatorname*{arg\,min}}
\newcommand{\inner}[2]{\langle #1,#2\rangle}
\newcommand{\norm}[1]{\|#1\|}
\begin{document}

\section*{Primal–Dual Hybrid Gradient (PDHG) Algorithm}

\section*{Mathematical Formulation}
Consider the convex–concave saddle‑point problem  
\[
\min_{x\in\mathbb R^{n}} \; g(x)+ \max_{y\in\mathbb R^{m}} \big\{ \langle Kx ,y\rangle - h^{*}(y)\big\},
\tag{1}
\]
where  

\begin{itemize}[noitemsep]
\item $g:\mathbb R^{n}\!\to\!\mathbb R\cup\{+\infty\}$ is proper, closed and convex,
\item $h^{*}:\mathbb R^{m}\!\to\!\mathbb R\cup\{+\infty\}$ is the convex conjugate of a proper closed convex function $h$,
\item $K\in\mathbb R^{m\times n}$ is a linear operator.
\end{itemize}
Denote a saddle point by $(x^{\star},y^{\star})$ satisfying
\[
g(x^{\star})+ \langle Kx^{\star},y\rangle - h^{*}(y)
\;\le\;
g(x^{\star})+ \langle Kx^{\star},y^{\star}\rangle - h^{*}(y^{\star})
\;\le\;
g(x)+ \langle Kx ,y^{\star}\rangle - h^{*}(y^{\star})
\quad\forall x,y .
\]

The PDHG iterates are defined for step sizes $\tau>0$, $\sigma>0$ as  

\[
\boxed{
\begin{aligned}
y_{k+1} &:= \operatorname*{prox}_{\sigma h^{*}}\!\bigl(y_{k}+ \sigma K\bar x_{k}\bigr), \tag{2a}\\[4pt]
x_{k+1} &:= \operatorname*{prox}_{\tau g}\!\bigl(x_{k}- \tau K^{\!\top} y_{k+1}\bigr), \tag{2b}\\[4pt]
\bar x_{k+1} &:= 2x_{k+1}-x_{k}. \tag{2c}
\end{aligned}}
\]

The extrapolation $\bar x_{k}=2x_{k}-x_{k-1}$ is the hallmark of the hybrid gradient scheme.

\section*{Pseudocode}
\begin{verbatim}
Input:   τ>0, σ>0, K, prox_{τ g}, prox_{σ h*}
Initialize: x0 ∈ R^n, y0 ∈ R^m,  x̄0 := x0
for k = 0,1,2,… do
    y_{k+1} = prox_{σ h*}( y_k + σ * K * x̄_k )
    x_{k+1} = prox_{τ g}( x_k - τ * K^T * y_{k+1} )
    x̄_{k+1} = 2 * x_{k+1} - x_k
end for
\end{verbatim}

Termination can be based on a maximum number of iterations $T$ or on the
primal–dual gap falling below a tolerance.

\section*{Dry Run Example}
We illustrate the algorithm on the scalar problem  

\[
\min_{x\in\mathbb R}\;(x-3)^{2},
\tag{3}
\]
which can be written in the form (1) with  

\[
g(x)= (x-3)^{2},\qquad h^{*}(y)=0,\qquad K=1.
\]

Since $h^{*}\equiv0$, its proximal operator is the identity:
$\operatorname{prox}_{\sigma h^{*}}(u)=u$.
The proximal of $g$ with step $\tau$ is
\[
\prox_{\tau g}(v)=\argmin_{x}\Bigl\{(x-3)^{2}+ \frac{1}{2\tau}(x-v)^{2}\Bigr\}
          =\frac{v+2\tau\cdot 3}{1+2\tau}.
\tag{4}
\]

Choose $\tau=\sigma=0.1$, $x_{0}=0$, $y_{0}=0$, and initialise $\bar x_{0}=x_{0}=0$.

\begin{center}
\begin{tabular}{c|c|c|c|c}
$k$ & $y_{k+1}$ & $x_{k+1}$ & $\bar x_{k+1}=2x_{k+1}-x_{k}$ & $f(x_{k+1})$\\\hline
0 &
$y_{1}=y_{0}+0.1\bar x_{0}=0$ &
$x_{1}= \frac{0+2\cdot0.1\cdot3}{1+2\cdot0.1}= \frac{0.6}{1.2}=0.5$ &
$\bar x_{1}=2\cdot0.5-0=1.0$ &
$(0.5-3)^{2}=6.25$\\\hline
1 &
$y_{2}=y_{1}+0.1\bar x_{1}=0+0.1\cdot1=0.1$ &
$x_{2}= \frac{0.5-0.1\cdot0.1+2\cdot0.1\cdot3}{1+2\cdot0.1}
      =\frac{0.5-0.01+0.6}{1.2}= \frac{1.09}{1.2}\approx0.9083$ &
$\bar x_{2}=2\cdot0.9083-0.5\approx1.3167$ &
$(0.9083-3)^{2}\approx4.380$\\\hline
2 &
$y_{3}=y_{2}+0.1\bar x_{2}=0.1+0.1\cdot1.3167=0.2317$ &
$x_{3}= \frac{0.9083-0.1\cdot0.2317+0.6}{1.2}
      =\frac{1.5083-0.0232}{1.2}\approx1.2376$ &
$\bar x_{3}=2\cdot1.2376-0.9083\approx1.5669$ &
$(1.2376-3)^{2}\approx3.099$\\
\end{tabular}
\end{center}
The objective value decreases from $9$ (at $x=0$) to $\approx3.10$ after three iterations, illustrating the descent of PDHG.

\section*{Assumptions}
\begin{assumption}[Convexity and Smoothness]
\label{ass:convex}
\begin{enumerate}
\item $g$ is convex and $L_{g}$‑Lipschitz differentiable, i.e.
\[
\|\nabla g(x)-\nabla g(x')\|\le L_{g}\|x-x'\|\quad\forall x,x'.
\]
\item $h^{*}$ is convex (no smoothness required; it may be non‑smooth).
\item $K$ has operator norm $\|K\|=L_{K}$.
\end{enumerate}
\end{assumption}

\begin{assumption}[Step‑size condition]
\label{ass:steps}
The primal and dual steps satisfy
\[
\tau\sigma L_{K}^{2}\le 1.
\tag{5}
\]
\end{assumption}

\section*{Main Convergence Theorem}
\begin{theorem}[Ergodic $\mathcal O(1/T)$ rate]\label{thm:rate}
Let $(x^{\star},y^{\star})$ be any saddle point of \eqref{1}.
Define the ergodic averages
\[
\bar x^{T}:=\frac{1}{T}\sum_{k=0}^{T-1}x_{k},\qquad
\bar y^{T}:=\frac{1}{T}\sum_{k=0}^{T-1}y_{k+1}.
\]
Under Assumptions \ref{ass:convex}–\ref{ass:steps}, the primal–dual gap satisfies
\[
\mathcal G(\bar x^{T},\bar y^{T})
:=\Bigl[g(\bar x^{T})+ \langle K\bar x^{T},\bar y^{T}\rangle- h^{*}(\bar y^{T})\Bigr]
 -\Bigl[g(x^{\star})+ \langle Kx^{\star},\bar y^{T}\rangle- h^{*}(\bar y^{T})\Bigr]
 \le
\frac{\|x_{0}-x^{\star}\|^{2}}{2\tau T}
+\frac{\|y_{0}-y^{\star}\|^{2}}{2\sigma T}.
\tag{6}
\]
Consequently $\mathcal G(\bar x^{T},\bar y^{T})= \mathcal O(1/T)$.
\end{theorem}

\section*{Theoretical Properties}
\begin{itemize}
\item \textbf{Stability.} Condition \eqref{ass:steps} guarantees that the
Lyapunov function 
\[
\Phi_{k}:=\frac{1}{2\tau}\|x_{k}-x^{\star}\|^{2}
          +\frac{1}{2\sigma}\|y_{k}-y^{\star}\|^{2}
\]
is non‑increasing up to the gap term (see Lemma~\ref{lem:descent}).
\item \textbf{Hyper‑parameter sensitivity.} The product $\tau\sigma$ must not exceed $1/L_{K}^{2}$.
Larger $\tau$ (or $\sigma$) accelerates the decrease of the primal (dual) part but forces the other step size to shrink.
\item \textbf{Comparison.}
For $\sigma\to0$ the method reduces to a forward–backward (prox‑gradient) scheme with rate $\mathcal O(1/T)$.
For $\tau=\sigma=1/L_{K}$ the algorithm matches the Chambolle–Pock optimal choice, yielding the same bound as in Theorem~\ref{thm:rate}.
\end{itemize}

\section*{Preliminaries (Definitions and Lemmas)}
\begin{lemma}[Three‑point identity for proximal operators]\label{lem:prox}
For any convex $p$ and $\lambda>0$,
\[
\inner{u-\prox_{\lambda p}(u)}{z-\prox_{\lambda p}(u)}
\ge \lambda\bigl(p(z)-p(\prox_{\lambda p}(u))\bigr)
\qquad\forall z .
\tag{7}
\]
\end{lemma}
\begin{proof}
Standard property of the proximal mapping; see e.g. \cite{parikh2014prox}.
\end{proof}

\begin{lemma}[Descent of the Lyapunov function]\label{lem:descent}
Let $(x_{k},y_{k})$ be generated by \eqref{2}. Then for any saddle point $(x^{\star},y^{\star})$
\[
\Phi_{k+1}\le \Phi_{k}
 -\bigl[g(x_{k+1})-g(x^{\star})+ h^{*}(y_{k+1})-h^{*}(y^{\star})
      +\inner{Kx_{k+1},y_{k+1}}-\inner{Kx^{\star},y_{k+1}}\bigr].
\tag{8}
\]
\end{lemma}
\begin{proof}
Apply Lemma~\ref{lem:prox} to the updates \eqref{2a} and \eqref{2b} with
$z=x^{\star}$ and $z=y^{\star}$ respectively, then use the step‑size condition
\eqref{ass:steps} to bound the cross term $\inner{K(x_{k+1}-x_{k}),y_{k+1}-y_{k}}$.
\end{proof}

\section*{Main Proof (Step‑by‑Step Derivation)}
We prove Theorem~\ref{thm:rate} by telescoping the inequality \eqref{8}.

\begin{proof}
Let $\Phi_{k}$ be defined as in Lemma~\ref{lem:descent}.
\begin{enumerate}[label=[Step \arabic*],leftmargin=*]
\item From Lemma~\ref{lem:descent} we have for every $k\ge0$
\[
\Phi_{k+1}\le \Phi_{k}
 -\Bigl[g(x_{k+1})-g(x^{\star})
      +h^{*}(y_{k+1})-h^{*}(y^{\star})
      +\inner{Kx_{k+1},y_{k+1}}-\inner{Kx^{\star},y_{k+1}}\Bigr].
\tag{9}
\]
\item Rearranging \eqref{9} yields
\[
g(x_{k+1})-g(x^{\star})
+h^{*}(y_{k+1})-h^{*}(y^{\star})
+\inner{Kx_{k+1},y_{k+1}}-\inner{Kx^{\star},y_{k+1}}
\le \Phi_{k}-\Phi_{k+1}.
\tag{10}
\]
\item Sum \eqref{10} over $k=0,\dots,T-1$:
\[
\sum_{k=0}^{T-1}\!\Bigl[g(x_{k+1})-g(x^{\star})
+h^{*}(y_{k+1})-h^{*}(y^{\star})
+\inner{Kx_{k+1},y_{k+1}}-\inner{Kx^{\star},y_{k+1}}\Bigr]
\le \Phi_{0}-\Phi_{T}.
\tag{11}
\]
Since $\Phi_{T}\ge0$, the right–hand side is bounded by $\Phi_{0}$.
\item Divide both sides of \eqref{11} by $T$ and use convexity of $g$ and $h^{*}$:
\[
g(\bar x^{T})-g(x^{\star})
+h^{*}(\bar y^{T})-h^{*}(y^{\star})
+\inner{K\bar x^{T},\bar y^{T}}-\inner{Kx^{\star},\bar y^{T}}
\le \frac{\Phi_{0}}{T}.
\tag{12}
\]
The inequality follows from Jensen’s inequality applied to each convex term.
\item By definition of the primal–dual gap $\mathcal G(\bar x^{T},\bar y^{T})$ in \eqref{6},
the left–hand side of \eqref{12} equals $\mathcal G(\bar x^{T},\bar y^{T})$.
Thus
\[
\mathcal G(\bar x^{T},\bar y^{T})\le \frac{\Phi_{0}}{T}
   =\frac{1}{2\tau T}\|x_{0}-x^{\star}\|^{2}
    +\frac{1}{2\sigma T}\|y_{0}-y^{\star}\|^{2}.
\tag{13}
\]
\item Inequality \eqref{13} is exactly the bound \eqref{6}, completing the proof.
\end{enumerate}
\end{proof}

\section*{Rate Derivation (With Constants)}
From \eqref{13} we read the explicit constants:
\[
C_{x}:=\frac{\|x_{0}-x^{\star}\|^{2}}{2\tau},
\qquad
C_{y}:=\frac{\|y_{0}-y^{\star}\|^{2}}{2\sigma},
\qquad
\mathcal G(\bar x^{T},\bar y^{T})\le \frac{C_{x}+C_{y}}{T}.
\]
If the primal and dual initial points are chosen at the origin and the
optimal solution has norm bounded by $R$, then $C_{x}+C_{y}\le
\frac{R^{2}}{2}\Bigl(\frac{1}{\tau}+\frac{1}{\sigma}\Bigr)$, yielding the familiar
$\mathcal O(1/T)$ rate.

\section*{Special Cases}
\begin{itemize}
\item \textbf{Pure gradient descent.}
If $h^{*}\equiv0$ and $K=I$, choosing $\sigma\to0$ forces $y_{k}=0$ for all $k$.
The update \eqref{2b} becomes
$x_{k+1}= \prox_{\tau g}(x_{k})$, i.e. a proximal gradient step with rate
$\mathcal O(1/T)$.
\item \textbf{Linear constraints.}
When $h^{*}$ is the indicator of a linear constraint $Ax=b$, the dual update
\eqref{2a} reduces to a gradient ascent on the Lagrange multiplier, and the
overall scheme coincides with the classical Arrow–Hurwicz method, again
enjoying the $\mathcal O(1/T)$ ergodic bound.
\end{itemize}

\begin{thebibliography}{9}
\bibitem{parikh2014prox}
N. Parikh and S. Boyd, ``Proximal Algorithms,'' \emph{Foundations and Trends in Optimization}, vol. 1, no. 3, pp. 127–239, 2014.
\end{thebibliography}

\end{document}